\documentclass[aps,prl,twocolumn,showpacs,superscriptaddress]{revtex4-2}

\usepackage{amsmath,amssymb,amsfonts}
\usepackage{graphicx}
\usepackage{booktabs}
\usepackage{longtable}
\usepackage{hyperref}
\usepackage{xcolor}
\usepackage{microtype}

\begin{document}

\title{W(3,3): Arithmetic Uniqueness, Ramanujan Tau Bridge, and Neutrino Mass}

\author{Wil Dahn}
\affiliation{Independent Researcher, Severna Park, Maryland, USA}
\email{wilcompute@users.noreply.github.com}

\date{\today}

\begin{abstract}
We study the Ramanujan bipartite graph $W(3,3)$, the unique member of the
Lubotzky--Phillips--Sarnak $W(3,q)$ family satisfying five simultaneous
arithmetic conditions. We prove that $W(3,3)$ is the unique graph in this
family where: (i) the zeta-function poles encode cyclotomic polynomial values
exactly; (ii) the spectral multiplicities satisfy $k+g = q^{q}$; (iii) the
spectral parameters reconstruct Ramanujan's tau function at $p=2,3$; (iv) the
eigenspace poles lie in a Heegner quadratic field; and (v) the first
post-barrier Euler factor of the Heegner CM curve equals the $W(3,3)$ Frobenius
eigenvalue. We establish that the tau-reconstruction horizon equals
$\Phi_6(3)=7$, coinciding with the conductor prime of the Heegner CM curve
$E_{-7}$. As an application, we derive a neutrino mass sum prediction
$\sum m_\nu = 0.101\,\text{eV}$ (NH) and a seesaw spectral cascade descending
through the $W(3,3)$ cyclotomic ladder at RH-neutrino scales
$10^{14.7\text{--}15.1}\,\text{GeV}$.
\end{abstract}

\pacs{02.10.Ox, 14.60.Pq, 11.30.Pb}

\maketitle

%% ============================================================
\section{The W(3,3) Ramanujan Graph}
\label{sec:graph}
%% ============================================================

\subsection{Definition and LPS Construction}

The Lubotzky--Phillips--Sarnak (LPS) family $W(3,q)$ consists of Cayley graphs
on $\mathrm{PSp}(4,\mathbb{F}_q)$ for prime powers $q$~\cite{lubotzky1988}.
For $q=3$ the resulting graph $W(3,3)$ is a $(k,k)$-biregular bipartite graph
with degree $k=12$, $v=40$ vertices per side, and all non-trivial eigenvalues
bounded by $2\sqrt{k-1}=2\sqrt{11}$, satisfying the Ramanujan
condition~\cite{lubotzky1988,sarnak1990}.

\subsection{Spectral Parameters}

Table~\ref{tab:params} summarises the key spectral and arithmetic parameters
of $W(3,3)$.

\begin{table}[h]
\centering
\caption{Spectral and arithmetic parameters of $W(3,3)$.}
\label{tab:params}
\begin{tabular}{lcc}
\toprule
\textbf{Parameter} & \textbf{Symbol} & \textbf{Value} \\
\midrule
Degree & $k$ & 12 \\
Vertices per side & $v$ & 40 \\
$f$-multiplicity & $f$ & 24 \\
$g$-multiplicity & $g$ & 15 \\
$r$-eigenvalue & $\lambda_r$ & 2 \\
$s$-eigenvalue & $\lambda_s$ & $-4$ \\
$\Phi_3(3)$ & $\Phi_3$ & 13 \\
$\Phi_4(3)$ & $\Phi_4$ & 10 \\
$\Phi_6(3)$ & $\Phi_6$ & 7 \\
$\mu = q+1$ & $\mu$ & 4 \\
$2k-1$ & --- & 23 \\
\bottomrule
\end{tabular}
\end{table}

\subsection{Ihara Zeta Function}

The Ihara zeta function of $W(3,3)$ factors as~\cite{ihara1966,bass1992}
\begin{align}
  p_1(u) &= 1 - 2u + 11u^2 \quad (r\text{-eigenspace sector}), \\
  p_2(u) &= 1 + 4u + 11u^2 \quad (s\text{-eigenspace sector}),
\end{align}
with constant term $k-1=11=q^2+q-1$ shared by both sectors.  The poles of
$p_1$ lie in $\mathbb{Q}(\sqrt{-\Phi_4})$ and those of $p_2$ in
$\mathbb{Q}(\sqrt{-\Phi_6})$.

\subsection{The Parameter Ring}

We define the $W(3,3)$ parameter ring as
\begin{equation}
  \mathcal{R} = \mathbb{Z}[k,\,g,\,f,\,v,\,\Phi_3,\,\Phi_4,\,\Phi_6,\,2k-1]
\end{equation}
with generators $\{12,15,24,40,13,10,7,23\}$.

%% ============================================================
\section{The W(3,3) Uniqueness Theorem}
\label{sec:uniqueness}
%% ============================================================

\begin{theorem}[$W(3,3)$ Uniqueness]
$W(3,3)$ is the unique $W(3,q)$ Ramanujan graph satisfying all five conditions
\emph{C1}--\emph{C5} simultaneously.
\end{theorem}

\textbf{C1 (Zeta pole--cyclotomic calibration).}
Define the deficit
\begin{equation}
  \Delta(q) = \frac{\operatorname{Im}(\text{poles of }p_1)^2}{4} - \Phi_4(q)
            = -\frac{(q-3)^2}{4}.
\end{equation}
Then $\Delta(q)=0$ if and only if $q=3$.

\textbf{C2 ($k+g=q^q$).}
\begin{equation}
  k(q)+g(q) = q(q+1) + \frac{q(q^2+1)}{2} = q^q \iff q=3.
\end{equation}
This reduces to $q^2+2q+3 = 2q^{q-1}$, which has the unique positive integer
root $q=3$.

\textbf{C3 (Ramanujan tau coincidence).}
\begin{align}
  -f(q) &= \tau(2) = -24 \iff q=3, \\
  k(q)\cdot q\cdot\Phi_6(q) &= \tau(3) = 252 \iff q=3.
\end{align}

\textbf{C4 (Heegner field).}
The eigenspace poles lie in $\mathbb{Q}(\sqrt{-\Phi_6(q)})$.  For $q=3$ this
is $\mathbb{Q}(\sqrt{-7})$, a Heegner field with class number $h=1$.
At $q=4$, $\Phi_6(4)=13$ gives $h(\mathbb{Q}(\sqrt{-13}))=2$.

\textbf{C5 (Post-barrier Frobenius match).}
\begin{equation}
  a_{k(q)-1}(E_{-\Phi_6(q)}) = \lambda_s(q) = -(q+1).
\end{equation}
For $q=3$: $a_{11}(E_{-7}) = -4 = \lambda_s(3)$.  Fails at $q=2$:
$a_5(E_{-3})=0 \neq \lambda_s(2)=-3$.

%% ============================================================
\section{The Ramanujan Tau Bridge}
\label{sec:tau}
%% ============================================================

\subsection{Reconstruction of \texorpdfstring{$\tau$}{tau} at
            \texorpdfstring{$2^a\cdot 3^b$}{2a3b}}

From $\tau(2)=-f=-24$ and $\tau(3)=k\cdot q\cdot\Phi_6=252$ the Hecke
recursion
\begin{equation}
  \tau(p^{n+1}) = \tau(p)\,\tau(p^n) - p^{11}\,\tau(p^{n-1})
\end{equation}
and multiplicativity $\tau(mn)=\tau(m)\tau(n)$ for $\gcd(m,n)=1$ reconstruct
all values $\tau(2^a\cdot 3^b)$ from $W(3,3)$ alone~\cite{ramanujan1916}.

\subsection{\texorpdfstring{$\tau(5)$}{tau(5)} and
            \texorpdfstring{$\tau(7)$}{tau(7)} in the W(3,3) Ring}

\begin{align}
  \tau(5) &= 2q(q-2)\,\Phi_6\,(2k-1) = 4830, \\
  \tau(7) &= -(q+1)\,\Phi_6\,\Phi_3\,(2k-1) = -16744.
\end{align}

\subsection{Ramanujan Congruence mod 23}

For all $n>3$ with $\gcd(n,23)=1$,
\begin{equation}
  \tau(n) \equiv 0 \pmod{2k-1} = 0 \pmod{23},
\end{equation}
with exceptions only at the $W(3,3)$ primes $q=3$ and $\mu-1=3$.

\subsection{The \texorpdfstring{$\Phi_6$}{Phi6} Barrier and
            Conductor--Barrier Theorem}

Values $\tau(p)$ lie in $\mathcal{R}$ for $p\leq\Phi_6(3)=7$; $\tau(11)$
requires the external prime 149.  The conductor of $E_{-7}$ is
$N(E_{-7})=49=\Phi_6(3)^2$, so the conductor prime equals the barrier prime.
The first prime above the barrier is $k-1=11$, giving
$a_{11}(E_{-7})=-4=\lambda_s$: the $W(3,3)$ Frobenius is the first external
Euler factor.

%% ============================================================
\section{Neutrino Mass Prediction}
\label{sec:neutrino}
%% ============================================================

\subsection{The \texorpdfstring{$\mu_{\mathrm{eff}}^2$}{mueff2}
            Spectral Parameter}

Define the effective spectral hierarchy parameter
\begin{equation}
  \mu_{\mathrm{eff}}^2(m) =
    -\frac{\log s^*(m)}{\log \Phi_4},
  \qquad
  s^*(m) = \frac{\text{geom.~mean}}{\text{max}},
\end{equation}
ranging from $0$ (quasi-degenerate, QD) to $\infty$ (fully hierarchical).

\subsection{W(3,3) Fixed-Point Candidates}

Table~\ref{tab:candidates} lists the four leading candidates ordered by
Bayesian posterior.

\begin{table}[h]
\centering
\caption{$W(3,3)$ neutrino mass candidates (CCCLV posterior).}
\label{tab:candidates}
\begin{tabular}{lccc}
\toprule
\textbf{Ordering} & \boldmath$\mu_{\mathrm{eff}}^2$ & \boldmath$\sum m_\nu$ (eV)
  & \textbf{Posterior} \\
\midrule
NH & $1/\mu = 1/4$ & 0.101 & 0.418 \\
IH & $1/\mu = 1/4$ & 0.110 & 0.285 \\
IH & $1/6$         & 0.122 & 0.154 \\
NH & $1/6$         & 0.128 & 0.106 \\
\bottomrule
\end{tabular}
\end{table}

\subsection{DESI Model Dependence}

Table~\ref{tab:desi} shows the $95\%$ upper limit on $\sum m_\nu$ from DESI DR1
under three dark energy models~\cite{desi2024}.

\begin{table}[h]
\centering
\caption{DESI DR1 $95\%$ upper limits and NH/$\mu_{\mathrm{eff}}^2=1/4$ status.}
\label{tab:desi}
\begin{tabular}{lcc}
\toprule
\textbf{DE Model} & \textbf{95\% UL (eV)} & \textbf{NH/1/4 status} \\
\midrule
$\Lambda$CDM   & 0.072 & Excluded \\
$w_0$CDM       & 0.113 & Allowed \\
$w_0 w_a$CDM   & 0.173 & Allowed \\
\bottomrule
\end{tabular}
\end{table}

\noindent DESI DR1 prefers $w_0 w_a$CDM at $\sim 2\sigma$ over $\Lambda$CDM.

\subsection{j-Tower CM Structure}

\begin{align}
  j(d=-4) &= k^3 = 1728, \\
  j(d=-7) &= -g^3 = -3375, \\
  j(d=-8) &= (v/2)^3 = 8000, \\
  j(d=-11) &= -2^{g} = -32768.
\end{align}

%% ============================================================
\section{Seesaw Cascade and the GUT Connection}
\label{sec:seesaw}
%% ============================================================

\subsection{Type-I Seesaw}

For Yukawa coupling $y_D=1$ and $v_{\mathrm{EW}}=246\,\text{GeV}$,
\begin{equation}
  M_R = \frac{(y_D\,v_{\mathrm{EW}})^2}{m_\nu}
      \sim 5.5\times10^{14}\text{--}1.4\times10^{15}\,\text{GeV},
\end{equation}
placing all three $M_i$ inside the leptogenesis window
$[10^9,10^{15}]\,\text{GeV}$ and below $M_{\mathrm{GUT}}$~\cite{buchmuller2005}.

\subsection{Seesaw Spectral Cascade}

Table~\ref{tab:cascade} shows the spectral RG descent through the $W(3,3)$
ladder.

\begin{table}[h]
\centering
\caption{Seesaw spectral cascade $T^0\to T^3$.}
\label{tab:cascade}
\begin{tabular}{llll}
\toprule
\textbf{Step} & \boldmath$\mu_{\mathrm{eff}}^2$
  & \textbf{Nearest W(3,3)} & \textbf{Scale} \\
\midrule
$T^0$ & $1/4$   & $1/\mu$       & LH neutrinos \\
$T^1$ & $0.140$ & $1/\Phi_6=1/7$ & RH neutrinos \\
$T^2$ & $0.075$ & $1/\Phi_3=1/13$ & (2nd seesaw) \\
$T^3$ & $0.040$ & $1/(2k-1)=1/23$ & (3rd seesaw) \\
\bottomrule
\end{tabular}
\end{table}

\subsection{Cascade Fixed Point}

The cascade map $T$ satisfies $T(\mu^*)=\mu^*$ uniquely at $\mu^*=0$ (QD
limit).  The convergence ratio
\begin{equation}
  \frac{T^{(n+1)}}{T^{(n)}} \to 0.5225 \approx \sqrt{\frac{\Phi_4}{\Phi_6^2}}
\end{equation}
confirms a spectral RG descent with no cyclic orbit.

%% ============================================================
\section{RG Perturbation Theory and
         \texorpdfstring{$\mu_{\mathrm{eff}}^2$}{mueff2} Precision Probe}
\label{sec:rg}
%% ============================================================

Running $\mu_{\mathrm{eff}}^2$ from the $W(3,3)$ fixed point at scale
$\Lambda_0 \sim M_R$ to low energy via one-loop SM RG equations~\cite{pdg2024}
yields a per-mille shift
\begin{equation}
  \delta\mu_{\mathrm{eff}}^2 \simeq
    \frac{\alpha_s(M_Z)}{\pi}\,\ln\frac{\Lambda_0}{M_Z} \lesssim 0.003,
\end{equation}
which is within reach of next-generation neutrino experiments (JUNO, DESI DR2,
KATRIN)~\cite{juno2022,katrin2022}.  A $5\sigma$ detection of
$\sum m_\nu = 0.101\,\text{eV}$ at DESI DR2 would directly probe the
$W(3,3)$ fixed-point structure.

%% ============================================================
\section*{Acknowledgements}
%% ============================================================

The author thanks the open-source Python and SageMath communities.  All
numerical results are reproducible from the companion code at
\url{https://github.com/wilcompute/W33-Theory}.

%% ============================================================
\bibliography{references}
%% ============================================================

\end{document}
